\section{Entropy versus a truncated inverse-gap tail}\label{sec:additive}

This section defines a scalar characterization that separates the part of
the law handled by rare-event scale entropy from a part handled by the
ordinary inverse-gap cost and a part left unresolved at the requested
accuracy. We prove its observable upper bound here and its matching lower
bound in Section~\ref{sec:matching-lower}. The decomposition is analytical:
the algorithm never observes a coin's component label.

\subsection{A homogeneous entropy and a selection functional}

For a finite submeasure $\sigma\le\mu$, use the canonical bands
$I_j=(b_j/2,b_j]$, $b_j=2^{-(j+1)}$, and put
\[
 a_j(\sigma)=\int_{I_j}d^{-2}\,\sigma(dd),\qquad
 A_\sigma=\sum_j a_j(\sigma).
\]
For a finite nonnegative vector $a$, define
\[
 \mathcal G(a)=A+\sum_j a_j\log(A/a_j),\qquad A=\sum_j a_j,
 \qquad \mathcal G(0)=0.
\]
Zero coordinates contribute zero. The functional $\mathcal G$ is concave, coordinatewise increasing,
homogeneous of degree one, and superadditive. Indeed, on positive
coordinates its derivative is $1+\log(A/a_j)$ and its Hessian is
$A^{-1}\mathbf1\mathbf1^\top-\operatorname{diag}(a_j^{-1})$,
which is negative semidefinite by Cauchy--Schwarz. Continuity gives the
boundary cases, and concavity plus homogeneity gives superadditivity. Define
\begin{equation}\label{eq:submeasure-entropy}
 \Gamma_\sigma=\sup_J\mathcal G(a_0(\sigma),\ldots,a_J(\sigma)).
\end{equation}
Thus $\Gamma_0=0$, and, when $0<A_\sigma<\infty$,
$\Gamma_\sigma=A_\sigma[1+\Ent((a_j/A_\sigma)_j)]$, with extended-value
interpretation. Also write
$H_\tau=\int d^{-1}\,\tau(dd)$ for any finite submeasure $\tau$.

Let $h=B^{-1/2}$ and define
\begin{equation}\label{eq:scalar-selection}
 \boxed{\quad
 \mathcal K_{\mu,\rho}(B)=
 \min\left\{\rho,\inf_{0\le\sigma\le\mu}
 \left[h\sqrt{\rho\Gamma_\sigma}
       +\int\min\{1,h/d\}\,d(\mu-\sigma)\right]\right\}.
 \quad}
\end{equation}
Equivalently,
\begin{equation}\label{eq:three-piece-functional}
 \mathcal K_{\mu,\rho}(B)=
 \min\left\{\rho,\inf_{\sigma+\tau+\zeta=\mu}
 \left[\zeta((0,1/2])
       +\frac{\sqrt{\rho\Gamma_\sigma}+H_\tau}{\sqrt B}\right]\right\}.
\end{equation}
Indeed, at fixed $\sigma$, assign each remaining magnitude to the cheaper
of the costs $1$ and $h/d$.

\begin{lemma}[Hard assignments suffice]\label{lem:hard-selection}
The infimum in \eqref{eq:scalar-selection} is unchanged if
$\sigma=\mu|S$ for a measurable set $S$. In particular, for a finite
atomic law a minimizing assignment can put each atom entirely in one of
the three pieces in \eqref{eq:three-piece-functional}.
\end{lemma}

\begin{proof}
Write $f=d\sigma/d\mu\in[0,1]$ and $S_t=\{f>t\}$, so
$\sigma=\int_0^1\mu|S_t\,dt$. Concavity of the square root of every
finite-fragment entropy gives
\[
 \sqrt{\Gamma_\sigma}\ge
       \int_0^1\sqrt{\Gamma_{\mu|S_t}}\,dt
\]
after monotone convergence in the fragment size. For a finite-objective
input the relevant entropies are finite for $t>0$, since
$\mu|S_t\le\sigma/t$ and
$\Gamma_{\mu|S_t}\le\Gamma_\sigma/t$. The remainder integral is linear,
so the objective at $\sigma$ is at least the average of the hard-set
objectives. Their infimum is therefore no larger. The reverse inequality
is immediate. For a finite law the minimum is attained among the finitely
many assignments. This is a selection problem, not a claim of convex
minimization over its assignments.
\end{proof}

The definition immediately implies
\[
 \mathcal K_{\mu,\rho}(B)\le
 \min\{\rho,\omega_\mu(B^{-1/2}),
                  \sqrt{\rho\Gamma_\mu/B}\},
 \qquad
 \omega_\mu(h)=\int\min\{1,h/d\}\,\mu(dd).
\]
Infinite quantities are allowed in this display.

\begin{theorem}[An entropy--tail law for the full joint problem]
\label{thm:entropy-tail-law}
There are universal constants $0<c<C<\infty$ such that, for every
probability measure $\mu$ on $(0,1/2]$, every $0<\rho\le1/4$, and every
$B>0$,
\begin{equation}\label{eq:entropy-tail-law}
 \boxed{\qquad
 c\,\mathcal K_{\mu,\rho}(B)\le R_{\mu,\rho}(B)
             \le C\,\mathcal K_{\mu,\rho}(B).
 \qquad}
\end{equation}
The functional has the equivalent selection and three-piece forms
\eqref{eq:scalar-selection}--\eqref{eq:three-piece-functional}, and hard
measurable assignments suffice. No density, regular variation, finite
inverse moment, or finite full-law entropy is required.
\end{theorem}

The upper half is Theorem~\ref{thm:three-piece-upper}; the lower half is
Theorem~\ref{thm:entropy-tail-lower}. The comparison constants in
\eqref{eq:entropy-tail-law} are universal but conservative. They are not
the factor sixteen proved for the finite convex program; these are
different representations and different comparison bounds.
The proof also gives the explicit conservative constants in
\eqref{eq:explicit-main-constants}.

\begin{theorem}[Three-piece observable upper bound]\label{thm:three-piece-upper}
There is a universal constant $C$ such that, for every magnitude law on
$(0,1/2]$, every $0<\rho\le1/4$, and every $B>0$,
\begin{equation}\label{eq:three-piece-upper}
 R_{\mu,\rho}(B)\le C\mathcal K_{\mu,\rho}(B).
\end{equation}
No component label or magnitude is revealed to the algorithm.
\end{theorem}

\subsection{An uncapped reference bank with a capped observable race}

Fix a decomposition $\mu=\sigma+\tau+\zeta$, put
$z=\zeta((0,1/2])$, and assume $\Gamma_\sigma,H_\tau<\infty$.
It suffices to treat $z\le\rho/128$: otherwise the constant $\rho/2$
has risk at most $64z$. Write $A=A_\sigma$, $H=H_\tau$.
If $\sigma\ne0$, use all its occupied bands, with weights
$\pi_j=a_j(\sigma)/A$ and references $b_j/2$ and $b_j/4$.
For an internal allowance $u<\rho$, set $\epsilon=u/\rho$.
Stop negative if a negative reference likelihood against fair reaches
$1/(\epsilon\pi_j)$; stop positive if
\begin{equation}\label{eq:split-bank-ratio}
 \frac{\sum_j\pi_jP_{1/2+b_j/4,n}}
      {\int P_{1/2-d,n}\,\mu(dd)}\ge1/u.
\end{equation}
Give negative crossings priority. A zero denominator triggers a positive
stop. The split $\sigma+\tau+\zeta=\mu$ specifies the references and the
work analysis; membership in these submeasures is never observed.
The known all-negative mixture in the denominator is the full law
$\mu$, regardless of the decomposition.

By Lemma~\ref{lem:uf-mixture-bank}, a wrong negative crossing has
conditional positive-sign probability at most $\epsilon$, and a positive
crossing has all-negative-mixture probability at most $u$. Thus the total
probability of a wrong bank stop is at most $2u$ over $\mathcal Q_\rho$.
Termination outside $\sigma$ is not assumed.

Race this bank against the ordinary classifier of Lemma~\ref{lem:uf-ordinary}, using
alternating observations of the same physical coin. At an ordinary cap
$M$, its conditional work certificate is
\[
 o_M(d)=\min\{M,17d^{-2}[1+\log_+(Md^2)]\}.
\]
For any finite measure $\xi$, put
\[
 \eta_\xi(M)=\int\min\{1,(Md^2)^{-2}\}\,\xi(dd),\qquad
 O_\xi(M)=\int o_M(d)\,\xi(dd).
\]
The ordinary error under $\mu$ is at most $\eta_\mu(M)$.
Stop as soon as either component stops and return that output.
On completed counterfactual streams, the physical work is at most
\[
 2\min\{T_{\rm bank},T_{\rm ordinary}\}\le2M.
\]
A wrong output implies either a wrong bank stop or an ordinary error.
When $\sigma=0$, omit the bank and use the ordinary classifier alone.

On a $\sigma$-almost-every gap $d$ in band $j$, the matching negative
reference has drift at least $4d^2/3$ and log overshoot below one.
Its expected crossing time is at most
\begin{equation}\label{eq:split-negative-work}
 \frac3{4d^2}[1+\log(1/\epsilon)+\log(1/\pi_j)].
\end{equation}
On a positive gap, the event $\widehat p_n-1/2\ge d/2$ gives
\begin{equation}\label{eq:split-null-laplace}
 \frac{P_{-,\mu,n}}{P_{{\rm f},n}}
 \le\int e^{-2nx^2}\,\mu(dx)
 \le z+\frac A{2en}+\frac H{2\sqrt{en}}.
\end{equation}
The last two bounds use, respectively,
$\sup_xx^2e^{-2nx^2}=1/(2en)$ and
$\sup_xxe^{-2nx^2}=1/(2\sqrt{en})$.
For $n\ge d^{-2}$ the denominator is at most $z+Ad^2+Hd$,
while the matching positive reference contributes at least
$\pi_je^{nd^2/3}$ relative to fair. Equation~\eqref{eq:uf-bank-positive} therefore bounds positive expected bank work by
\begin{equation}\label{eq:split-positive-work}
 \frac6{d^2}\left[
 1+\log(1/\epsilon)+\log(1/\pi_j)
 +\log_+\frac{z+Ad^2+Hd}{\rho}\right].
\end{equation}
These arguments prove finite conditional bank work on $\sigma$.
On $\tau+\zeta$, the ordinary cap guarantees termination of the race.

Since $z/\rho\le1$,
\[
 \log_+\frac{z+Ad^2+Hd}{\rho}
 \le\log3+\log_+(Ad^2/\rho)+\log_+(Hd/\rho).
\]
Using $\sup_{x\ge1}\log x/x=1/e$ and
$\sup_{x\ge1}\log x/x^2=1/(2e)$ gives
\begin{equation}\label{eq:split-surcharge}
 \int q(d)d^{-2}\log_+\frac{z+Ad^2+Hd}{\rho}\,\sigma(dd)
 \le(\log3+1/e)A+\frac{H^2}{2e\rho}.
\end{equation}
Here $\int q\,d\sigma\le\rho$ and $q\le1$ follow from the original
global rarity class; no separate conditional rarity assumption is used.
Integrating \eqref{eq:split-negative-work} and
\eqref{eq:split-positive-work} therefore bounds the bank work on $\sigma$ by
\begin{equation}\label{eq:split-bank-work}
 6\Gamma_\sigma+6A\log(1/\epsilon)
       +6(\log3+1/e)A+\frac{3H^2}{e\rho}.
\end{equation}

\subsection{A resolved-law cap schedule}

Let $\nu=\sigma+\tau$ and set
\[
 \delta_\ell=\rho16^{-(\ell+1)},\qquad
 u_\ell=\delta_\ell/4,\qquad \epsilon_\ell=u_\ell/\rho.
\]
Choose $M_\ell$ as the least power of four, at least four, such that
$\eta_\nu(M_\ell)\le\delta_\ell/8$.
Every chosen cap has a predecessor because $\eta_\nu(4)=1-z$.
The inequality $\eta_\nu(4M)\ge\eta_\nu(M)/16$ implies
\begin{equation}\label{eq:split-cap-selection}
 \delta_\ell/128<\eta_\nu(M_\ell)\le\delta_\ell/8.
\end{equation}
Thus the selected powers of four are distinct and increasing.

If $z>0$, let $J$ be the largest integer with $\delta_J\ge8z$.
Then
\begin{equation}\label{eq:split-terminal-level}
 8z\le\delta_J<128z.
\end{equation}
For $\ell\le J$,
$\eta_\mu(M_\ell)\le\eta_\nu(M_\ell)+z\le\delta_\ell/4=u_\ell$.
The raced classifier has total error at most $3u_\ell\le\delta_\ell$
and false-positive probability at most $2u_\ell\le\delta_\ell$.
When $z=0$, these conclusions hold at every level.

Since $\log(1/\epsilon_\ell)=\log64+\ell\log16$ and
$A\le\Gamma_\sigma$, \eqref{eq:split-bank-work} is at most
$40(\ell+1)\Gamma_\sigma+2H^2/\rho$.
A uniform raced-work certificate is consequently
\begin{equation}\label{eq:split-level-work}
 C_\ell=80(\ell+1)\Gamma_\sigma+4H^2/\rho
                  +2O_\tau(M_\ell)+2zM_\ell.
\end{equation}
It remains a valid, possibly loose, bound when $\sigma=0$ and the bank is
omitted.

The binary level-zero output has second moment at most
$\rho+\delta_0=17\delta_0$. An adjacent difference, run on successive
unused blocks of the same coin, has second moment at most
$\delta_\ell+\delta_{\ell-1}=17\delta_\ell$.
Set $v_\ell=17\delta_\ell$, $k_0=C_0$,
$k_\ell=C_\ell+C_{\ell-1}$ for $\ell\ge1$, and
$S_J=\sum_{\ell\le J}\sqrt{v_\ell k_\ell}$.
Independent-Poisson allocation across these corrections
(Lemma~\ref{lem:uf-poisson})
gives RMS at most
\begin{equation}\label{eq:split-poisson}
 \sqrt{\delta_J^2+S_J^2/B}.
\end{equation}
The seed may choose zero replications, and the Poisson means may be
below one. Every first and subsequent observation is charged.
Every such finite-level procedure has finite deterministic local caps.
The reference bank may be countable; its stopping events are measurable
countable unions and its positive reference is a probability mixture.

\subsection{Controlling mixed work without an extra logarithm}

For every finite $\tau$, the following bounds hold:
\begin{equation}\label{eq:split-elementary-kernels}
 O_\tau(M)\le21H\sqrt M,\qquad
 \eta_\sigma(M)\le A/M,\qquad
 \eta_\tau(M)\le H/\sqrt M.
\end{equation}
The latter two are pointwise inequalities before integration.
For the first, put $x=\sqrt M\,d$. The ratio of $o_M(d)$ to
$\sqrt M/d$ is at most one for $x\le1$, and at most
$17(1+2\log x)/x\le34/\sqrt e<21$ for $x\ge1$.
Lemma~\ref{lem:uf-kernel} supplies the homogeneous ordinary kernel bound:
\begin{equation}\label{eq:split-ordinary-kernel}
 \sum_{k\ge1}\sqrt{\eta_\tau(4^k)O_\tau(4^k)}\le 208H.
\end{equation}
Its proof applies without normalization of $\tau$: dyadic binning yields
the summable kernels $2^{k-j}$ for $j\ge k$, and
$2^{3(j-k)}$, $2^{j-k}(1+k-j)$ for $j<k$.
Apply $\sqrt{xy}\le(x+y)/2$ and sum these kernels.

Split levels according to which term of
$\eta_\nu=\eta_\sigma+\eta_\tau$ is at least half the total.
On levels where $\eta_\tau$ dominates,
$\delta_\ell\le256\eta_\tau(M_\ell)$, so distinctness of the caps and
\eqref{eq:split-ordinary-kernel} give a total
$\sum\sqrt{\delta_\ell O_\tau(M_\ell)}\le3328H$.
On the other levels,
\[
 \delta_\ell\le256\eta_\sigma(M_\ell)\le256A/M_\ell,
 \qquad M_\ell\le256A/\delta_\ell.
\]
Therefore
\[
 \sqrt{\delta_\ell O_\tau(M_\ell)}
 \le4\sqrt{21}\,H^{1/2}A^{1/4}\delta_\ell^{1/4}.
\]
The sum of $\delta_\ell^{1/4}$ over all levels is $\rho^{1/4}$.
The arithmetic--geometric mean inequality proves
\begin{equation}\label{eq:split-mixed-sum}
 \sum_{\ell\le J}\sqrt{\delta_\ell O_\tau(M_\ell)}
 \le C_{\rm mix}[H+\sqrt{\rho A}].
\end{equation}

The work on the unresolved mass is controlled separately.
On levels where $\eta_\sigma$ dominates,
\begin{equation}\label{eq:split-unresolved-sigma}
 \sum\sqrt{z\delta_\ell M_\ell}
 \le16(J+1)\sqrt{zA}\le\sqrt{2\rho A}.
\end{equation}
The last inequality uses $z\le\delta_J/8$ and
$(J+1)4^{-(J+1)}\le1/4$.
On levels where $\eta_\tau$ dominates,
$\sqrt{M_\ell}\le256H/\delta_\ell$, so
\begin{equation}\label{eq:split-unresolved-tau}
 \sum\sqrt{z\delta_\ell M_\ell}
 \le256H\sqrt z\sum_{\ell=0}^J\delta_\ell^{-1/2}
 \le\frac{1024}{3}H\sqrt{z/\delta_J}
 \le\frac{1024}{3\sqrt8}H.
\end{equation}
When $z=0$ these contributions vanish.

\begin{proof}[Completion of Theorem~\ref{thm:three-piece-upper}]
The shifted adjacent-level sum has a factor $1/4$, so
\[
 S_J\le\frac{5\sqrt{17}}4
                \sum_{\ell=0}^J\sqrt{\delta_\ell C_\ell}.
\]
Use \eqref{eq:split-level-work}, \eqref{eq:split-mixed-sum},
\eqref{eq:split-unresolved-sigma}, \eqref{eq:split-unresolved-tau},
and $A\le\Gamma_\sigma$. The elementary sums
$\sum_{\ell\ge0}4^{-\ell}\sqrt{\ell+1}\le16/9$ and
$\sum_{\ell\ge0}4^{-\ell}=4/3$ then give
\[
 S_J\le C_{\rm sum}[\sqrt{\rho\Gamma_\sigma}+H].
\]
For $z>0$, \eqref{eq:split-terminal-level} and
\eqref{eq:split-poisson} imply
$R(B)\le128z+C_{\rm sum}[\sqrt{\rho\Gamma_\sigma}+H]/\sqrt B$.
For $z=0$, the bound on $S_J$ holds at every finite terminal level.
Use \eqref{eq:split-poisson} and then $\delta_J\to0$; taking the infimum
over these finite constructions gives the asserted bound.
The constant estimator covers $z>\rho/128$.
Finally take the infimum over decompositions and compare with $\rho/2$.
All approximations are at fixed $(\mu,\rho,B)$. The comparison constant
is independent of the chosen decomposition and terminal level; no
attaining infinite-dimensional decomposition is required.
No infinite-algorithm interchange is required. The same infimum can be
approached with finite reference banks: finite $\Gamma_\sigma$ implies
$H_\sigma<\infty$, so truncate $\sigma$ to finitely many bands and move its
remainder to $\tau$. The entropy does not increase and the additional
inverse-gap cost tends to zero. Then use the finite-level construction.
\end{proof}
